% ASSERT_CONVENTION: natural_units=natural, metric_signature=mostly_minus, fourier_convention=physics, coupling_convention=alpha_s, generator_normalization=T(fund)=1/2, covariant_derivative_sign=D_mu=partial_mu-igA
%
% Lecture 8: Spectrum, Scales, and Outlook
% Part of: The Dualised Standard Model -- Part III Lecture Notes
%
% Conventions (Seiberg hep-th/9411149, unchanged from Lectures 1-7):
%   T(fund) = 1/2 per Weyl fermion
%   b_0 = (11N_c - 2N_f)/3 for N_f Dirac flavours
%   A(fund) = 1 (cubic anomaly coefficient per fundamental Weyl fermion)
%   alpha_s = g_s^2 / (4 pi)
%   R[Q] = (N_f - N_c)/N_f; R[W] = 2
%   N_f counts Dirac flavour pairs (each = 2 Weyl fermions)
%
% CONVENTIONS CARRIED FROM LECTURE 7:
%   Hypercharge: Y = T_R^3 + Q_V/6  (Cacciapaglia-Sannino 2407.17281 Eq. 6)
%   Epistemic: all non-SUSY duality claims carry conjectural qualifiers (P1)
%   Universal Yukawa: y (single coupling for all fermions)
%   VEV convention: v^2 = sum_{ij} (|H^u_{ij}|^2 + |H^d_{ij}|^2) = (246 GeV)^2/2
%
\documentclass[11pt,a4paper]{article}
% ASSERT_CONVENTION: natural_units=natural, metric_signature=mostly_minus, fourier_convention=physics, coupling_convention=alpha_s, generator_normalization=T(fund)=1/2, covariant_derivative_sign=D_mu=partial_mu-igA
%
% CONVENTIONS (Seiberg hep-th/9411149)
% Metric: (+,-,-,-) (mostly minus)
% T(fund) = 1/2 per Weyl fermion
% b_0 = (11N_c - 2N_f)/3 for N_f Dirac flavours
% D_mu = partial_mu - ig A_mu^a T^a
% A(fund) = 1 (cubic anomaly coefficient per fundamental Weyl fermion)
% alpha_s = g_s^2 / (4 pi)
% R[Q] = (N_f - N_c)/N_f; R[W] = 2
% N_f counts Dirac flavour pairs (each = 2 Weyl fermions)
%
% This file is \input{} by all lecture files.

%% ---------- Document class ----------
% (loaded by the wrapper; preamble only provides packages and macros)

%% ---------- Standard packages ----------
\usepackage{amsmath,amssymb,amsthm,mathtools}
\usepackage{hyperref}
\usepackage[capitalise,noabbrev]{cleveref}
\usepackage{enumitem}
\usepackage{booktabs}
\usepackage{array}
\usepackage{xcolor}
\usepackage{tikz}
\usetikzlibrary{arrows.meta,decorations.markings,positioning,calc}
\usepackage{fancybox}
\usepackage{tcolorbox}
\tcbuselibrary{skins,breakable}
\usepackage{slashed}              % Feynman slash notation
\usepackage{cite}                 % compressed citation lists

%% ---------- Hyperref configuration ----------
\hypersetup{
  colorlinks=true,
  linkcolor=blue!70!black,
  citecolor=green!50!black,
  urlcolor=blue!80!black,
  pdfauthor={Part III Lecture Notes},
  pdftitle={The Dualised Standard Model}
}

%% ---------- Theorem environments ----------
\theoremstyle{plain}
\newtheorem{theorem}{Theorem}[section]
\newtheorem{lemma}[theorem]{Lemma}
\newtheorem{proposition}[theorem]{Proposition}
\newtheorem{corollary}[theorem]{Corollary}

\theoremstyle{definition}
\newtheorem{definition}[theorem]{Definition}
\newtheorem{example}[theorem]{Example}
\newtheorem{exercise}{Exercise}[section]

\theoremstyle{remark}
\newtheorem{remark}[theorem]{Remark}

%% ---------- Physics macros: gauge groups and representations ----------
\newcommand{\Nc}{N_{\!c}}
\newcommand{\Nf}{N_{\!f}}
\newcommand{\SU}[1]{\mathrm{SU}(#1)}
\newcommand{\UU}[1]{\mathrm{U}(#1)}

% Representations
\newcommand{\fund}{\square}
\newcommand{\antifund}{\overline{\square}}
\newcommand{\adj}{\mathrm{adj}}
\newcommand{\antisym}{\wedge^{\!2}}        % antisymmetric rank-2
\newcommand{\sym}{\mathrm{S}_2}           % symmetric rank-2

%% ---------- Physics macros: group theory constants ----------
% Dynkin index: Tr(T^a T^b) = T(R) delta^{ab}
\newcommand{\Tfund}{\tfrac{1}{2}}          % T(fund) = 1/2 per Weyl fermion
\newcommand{\Tadj}{\Nc}                    % T(adj) = N_c
\newcommand{\Cfund}{\frac{\Nc^2 - 1}{2\Nc}}  % C_2(fund)
\newcommand{\Cadj}{\Nc}                    % C_2(adj)

% Cubic anomaly coefficient: A(R) = Tr_R[T^a {T^b, T^c}] with A(fund) = 1
\newcommand{\Afund}{1}

%% ---------- Physics macros: beta function ----------
\newcommand{\bzero}{b_0}
\newcommand{\bone}{b_1}
\newcommand{\alphas}{\alpha_s}
\newcommand{\gs}{g_s}

%% ---------- Physics macros: SUSY / R-charge ----------
\newcommand{\Rcharge}[1]{R[#1]}
\newcommand{\Wsup}{W}                      % superpotential

%% ---------- Physics macros: covariant derivative and field strength ----------
\newcommand{\Dmu}{D_\mu}
\newcommand{\Fmn}{F_{\mu\nu}}
\newcommand{\Ftilde}{\widetilde{F}_{\mu\nu}}

%% ---------- Physics macros: common abbreviations ----------
\newcommand{\ie}{\textit{i.e.}}
\newcommand{\eg}{\textit{e.g.}}
\newcommand{\cf}{\textit{cf.}}
\newcommand{\IR}{\ensuremath{\mathrm{IR}}}
\newcommand{\UV}{\ensuremath{\mathrm{UV}}}
\newcommand{\AF}{\ensuremath{\mathrm{AF}}}
\newcommand{\BZ}{\ensuremath{\mathrm{BZ}}}
\newcommand{\NSVZ}{\ensuremath{\mathrm{NSVZ}}}
\newcommand{\SQCD}{\ensuremath{\mathrm{SQCD}}}
\newcommand{\QCD}{\ensuremath{\mathrm{QCD}}}
\newcommand{\SM}{\ensuremath{\mathrm{SM}}}
\newcommand{\Tr}{\mathrm{Tr}}

%% ---------- Convention box macro ----------
\newtcolorbox{conventionbox}[1][]{%
  enhanced,
  colback=blue!5!white,
  colframe=blue!60!black,
  fonttitle=\bfseries,
  title={Conventions},
  #1
}

%% ---------- Comparison table style ----------
\newtcolorbox{comparisontable}[1][]{%
  enhanced,
  colback=orange!5!white,
  colframe=orange!70!black,
  fonttitle=\bfseries,
  title={Comparison},
  #1
}

%% ---------- Warning / important box ----------
\newtcolorbox{warningbox}[1][]{%
  enhanced,
  colback=red!5!white,
  colframe=red!70!black,
  fonttitle=\bfseries,
  title={Important},
  #1
}

%% ---------- Preview / forward-reference box ----------
\newtcolorbox{previewbox}[1][]{%
  enhanced,
  colback=green!5!white,
  colframe=green!50!black,
  fonttitle=\bfseries,
  title={Preview},
  #1
}

%% ---------- Page layout ----------
\usepackage[margin=2.5cm]{geometry}
\setlength{\parskip}{0.4em}
\setlength{\parindent}{0em}

%% ---------- Section formatting ----------
\usepackage{titlesec}
\titleformat{\section}{\Large\bfseries}{\thesection}{1em}{}
\titleformat{\subsection}{\large\bfseries}{\thesubsection}{1em}{}

%% ---------- End of preamble ----------


\title{\textbf{Lecture 8: Spectrum, Scales, and Outlook}}
\author{The Dualised Standard Model --- Part III Lecture Notes}
\date{Lent Term 2027}

\begin{document}
\maketitle

\tableofcontents
\newpage

%% ========================================================================
%%  CONVENTIONS
%% ========================================================================

\begin{conventionbox}[title={Conventions for Lecture 8}]
All conventions from Lectures~1--7 remain in force.  This lecture adds
no new conventions but makes extensive use of the Lecture~7 hypercharge
formula and the universal Yukawa coupling structure.

\renewcommand{\arraystretch}{1.3}
\begin{tabular}{@{}l l l@{}}
\toprule
\textbf{Quantity} & \textbf{Convention} & \textbf{Comment} \\
\midrule
Units & $\hbar = c = 1$ (natural) & \\
Metric & $\eta_{\mu\nu} = \mathrm{diag}(+1,-1,-1,-1)$ & mostly minus \\
Dynkin index & $\Tr(T^a T^b) = \Tfund\, \delta^{ab}$ & fundamental of $\SU{\Nc}$ \\
Hypercharge & $Y = T_R^3 + \tfrac{1}{6}\,Q_V$ & Lecture~7; Eq.~(6) of \cite{CacciapagliaSannino2024} \\
Universal Yukawa & $y$ (single coupling) & Eq.~(4) of \cite{CacciapagliaSannino2024} \\
VEV sum rule & $v^2 = \sum_{ij}(\langle H_{ij}^u\rangle^2 + \langle H_{ij}^d\rangle^2) = \tfrac{1}{2}(246\;\text{GeV})^2$ & input from $G_F$ \\
Epistemic & ``if the duality holds'' & all non-SUSY claims (P1) \\
\bottomrule
\end{tabular}

\medskip

\textbf{Parametric consistency framing.}  The numerical values
presented in this lecture (the top quark mass, the QCD scale, the Fermi
scale) are \textbf{parametric consistency checks} with $\mathcal{O}(1)$
free coefficients---they are \emph{not} first-principles predictions.
The dualised SM provides a framework in which these scales arise
\emph{naturally}, but the precise numerical values depend on free
parameters (the duality scale $\Lambda_D$ and the Planck-scale operator
coefficients $c^{abcd}$).
\end{conventionbox}


%% ========================================================================
%%  SECTION 1: Universal Yukawa and Scalar Democracy
%% ========================================================================
\section{Universal Yukawa and Scalar Democracy}
\label{sec:scalar-democracy}

In Lecture~7, we established that the magnetic Lagrangian contains a
\textbf{single universal Yukawa coupling} $y$ between the magnetic
quarks and the 18 Higgs doublets
(Eq.~\eqref{eq:yukawa-lecture7} below, reproduced from Lecture~7
Eq.~(12)).  This section develops the quantitative consequences of this
structure.

\subsection{From universal coupling to effective Yukawa matrices}

Recall from Lecture~7 that the magnetic Yukawa interaction takes the
form~\cite{CacciapagliaSannino2024}:
\begin{equation}\label{eq:yukawa-lecture7}
  y\, q\, \widetilde{q}\, \Phi_H
  \;=\; y \sum_{i,j}
  \Bigl(
    q_L^i\, u_R^j\, H_{ij}^u
    \;+\; q_L^i\, d_R^j\, H_{ij}^d
  \Bigr)\,,
\end{equation}
where $i,j = 1,2,3$ are generation indices and $y$ is the
\textbf{same coupling} for every generation pair.  This is a radical
departure from the Standard Model, where each fermion species has its
own independent Yukawa coupling $y_t$, $y_b$, $y_\tau$, $\ldots$ that
must be measured experimentally.

When the Higgs fields $H_{ij}^{u,d}$ acquire vacuum expectation values,
the fermion mass matrices are generated.  The \textbf{effective Yukawa
couplings} are~\cite{CacciapagliaSannino2024,HillMachado2019}:
\begin{equation}\label{eq:eff-yukawa-8}
  \boxed{%
  Y_{ij}^u \;=\; y\, \frac{\langle H_{ij}^u \rangle}{v}\,,
  \qquad
  Y_{ij}^d \;=\; y\, \frac{\langle H_{ij}^d \rangle}{v}\,,}
\end{equation}
(Eq.~(11) of~\cite{CacciapagliaSannino2024}), where
\begin{equation}\label{eq:vev-sum-8}
  v^2 \;=\; \sum_{i,j}
  \Bigl(
    \langle H_{ij}^u \rangle^2 + \langle H_{ij}^d \rangle^2
  \Bigr)
  \;=\; \frac{1}{2}\,(246\;\text{GeV})^2
\end{equation}
is the electroweak VEV, fixed by the measured Fermi constant $G_F$.

\begin{remark}[The VEV sum rule is an input, not a prediction]
  The value $v = 246/\sqrt{2}\;\text{GeV} \approx 174\;\text{GeV}$ is
  \emph{not} predicted by the dualised SM framework.  It is an
  experimental input---the Fermi constant $G_F = 1/(\sqrt{2}\, v^2)$
  is measured to high precision.  What the framework \emph{does} provide
  is a mechanism for distributing this total VEV hierarchically among
  the 18 Higgs doublets.
\end{remark}

\subsection{The scalar democracy mechanism}

The key insight, first articulated in this context by Hill, Machado,
Thomsen, and Turner~\cite{HillMachado2019}, is the \textbf{scalar
democracy mechanism}: the fermion mass hierarchy can be generated
entirely from a hierarchy in the Higgs VEVs, with no need for a
hierarchy in the Yukawa couplings.

In the Standard Model, the observed mass hierarchy
\begin{equation}\label{eq:mass-hierarchy}
  m_t \;\gg\; m_b \;\gg\; m_c \;\gg\; m_s \;\gg\; m_d \;\gg\; m_u
\end{equation}
is encoded in the Yukawa couplings: $y_t \approx 1$,
$y_b \approx 0.024$, $y_c \approx 0.007$, and so on, spanning five
orders of magnitude.  These couplings are free parameters in the SM
Lagrangian, and the origin of their hierarchy is one of the major
unsolved problems of particle physics---the \textbf{flavour puzzle}.

In the dualised SM, \emph{if the duality holds}, the situation is
fundamentally different.  Since the Yukawa coupling $y$ is universal,
the mass hierarchy must arise from the VEV hierarchy:
\begin{equation}\label{eq:mass-from-vev}
  \frac{m_i}{m_j}
  \;=\; \frac{Y_{ij}}{Y_{kl}}
  \;=\; \frac{\langle H_{ij} \rangle}{\langle H_{kl} \rangle}\,.
\end{equation}
The flavour puzzle is thus \emph{translated} from the Yukawa sector to
the scalar sector: instead of asking ``why are the Yukawa couplings
hierarchical?'', we ask ``why are the Higgs VEVs hierarchical?''  As we
will see in Section~\ref{sec:planck-operators}, the VEV hierarchy has a
natural origin in Planck-scale operators.

\subsection{Top quark Yukawa: $y_t \sim \mathcal{O}(1)$}

The dominant VEV is assigned to the $(3,3)$ up-type Higgs doublet:
$H_0 \equiv H_{33}^u$, which acquires
$\langle H_0 \rangle \approx v$.  (The $(3,3)$ label is a convention
reflecting that the top quark couples most strongly to this doublet.)
The top quark Yukawa coupling is then:
\begin{equation}\label{eq:top-yukawa}
  y_t \;=\; y \,\frac{\langle H_{33}^u \rangle}{v}
  \;\approx\; y \cdot 1
  \;=\; y\,.
\end{equation}
\emph{If the duality holds}, the observed large top Yukawa coupling
$y_t \approx 1$ is \textbf{natural}: it is simply the universal
magnetic Yukawa coupling itself.  The top mass is:
\begin{equation}\label{eq:top-mass}
  m_t \;\approx\; y\, v
  \;\sim\; \mathcal{O}(v)
  \;\sim\; 174\;\text{GeV}\,,
\end{equation}
which is parametrically consistent with the measured value
$m_t = 172.7 \pm 0.3\;\text{GeV}$.

\begin{warningbox}[title={Parametric consistency, not prediction}]
  The statement $m_t \sim \mathcal{O}(v)$ is a \textbf{parametric
  consistency check}, not a first-principles prediction.  The
  $\mathcal{O}(1)$ coefficient $y$ is a free parameter of the magnetic
  theory, determined by the strong dynamics at the duality scale
  $\Lambda_D$.  The dualised SM does \emph{not} predict
  $m_t = 172.7\;\text{GeV}$---it provides a framework where
  $m_t \sim v$ is the \emph{natural} scale, in contrast to the SM where
  $y_t \approx 1$ is an unexplained numerical coincidence.
\end{warningbox}

\subsection{Bottom-to-top mass ratio}

The bottom quark mass arises from the $(3,3)$ down-type Higgs VEV:
\begin{equation}\label{eq:bottom-yukawa}
  y_b \;=\; y\, \frac{\langle H_{33}^d \rangle}{v}\,,
  \qquad
  \frac{m_b}{m_t}
  \;=\; \frac{\langle H_{33}^d \rangle}{\langle H_{33}^u \rangle}\,.
\end{equation}
The measured ratio $m_b / m_t \approx 4.2 / 173 \approx 0.024$
requires:
\begin{equation}\label{eq:vev-ratio-bt}
  \langle H_{33}^d \rangle
  \;\sim\; 0.024\, v
  \;\sim\; 4\;\text{GeV}\,,
\end{equation}
which is the VEV of a ``second-tier'' Higgs doublet---much smaller than
the leading VEV $\langle H_0 \rangle \sim v$, but non-zero.  The
origin of this suppression is the subject of the next section.


%% ========================================================================
%%  SECTION 2: Planck-Scale Flavour Operators
%% ========================================================================
\section{Planck-Scale Flavour Operators}
\label{sec:planck-operators}

The VEV hierarchy required for the fermion mass pattern of
Section~\ref{sec:scalar-democracy} does not arise spontaneously from
the tree-level scalar potential of the 18 Higgs doublets.  Instead,
\emph{if the duality holds}, it is generated by Planck-scale
operators---higher-dimensional interactions suppressed by the Planck
mass $M_{\text{Pl}} \approx 1.2 \times 10^{19}\;\text{GeV}$.

\subsection{Four-fermion operators in the electric theory}

In the electric theory, the flavour symmetries of the fermion sector
are broken by gravity.  Since gravity couples to all forms of energy and
does not respect global symmetries, we expect operators of the
form~\cite{CacciapagliaSannino2024}:
\begin{equation}\label{eq:planck-operator}
  \boxed{%
  \mathcal{L}_{\text{Planck}}
  \;\supset\;
  \frac{c^{abcd}}{M_{\text{Pl}}^2}\,
  (Q^a \widetilde{Q}^b)(Q^c \widetilde{Q}^d)^\dagger\,,}
\end{equation}
(Eq.~(21) of~\cite{CacciapagliaSannino2024}), where $a,b,c,d$ are
flavour indices running over all $\Nf = 6$ flavours, and the
\textbf{coefficients $c^{abcd}$ are $\mathcal{O}(1)$ numbers} whose
precise values are unknown.

These operators are dimension-6 ($[Q]^4 / [M_{\text{Pl}}]^2$) and
are suppressed by two powers of the Planck mass.  They are the
\emph{leading} source of flavour violation in the electric theory:
all lower-dimensional operators respect the full
$\SU{6}_L \times \SU{6}_R$ flavour symmetry.

\subsection{Scalar mass terms from Planck operators}

At the duality scale $\Lambda_D$, the electric four-fermion operators
of Eq.~\eqref{eq:planck-operator} map to scalar mass terms in the
magnetic theory.  Since the mes-Higgs field
$\Phi_H = \langle \widetilde{Q}\lambda\lambda Q \rangle$ is a
composite of the electric fermions, the Planck operators generate both
\textbf{off-diagonal} mixing masses and \textbf{diagonal} masses for
the Higgs doublets:
\begin{itemize}
  \item \textbf{Off-diagonal:} Mixing mass terms $\mu_{0i}^2$ that
    connect the leading Higgs doublet $H_0 \equiv H_{33}^u$ to the
    remaining 17 ``dormant'' Higgs doublets $H_i$
    ($i = 1, \ldots, 17$).
  \item \textbf{Diagonal:} Large positive mass-squared terms $M_{ij}^2$
    for the dormant Higgs doublets, keeping them heavy and preventing
    them from acquiring large VEVs on their own.
\end{itemize}
The scale of these mass terms is set by dimensional
analysis~\cite{CacciapagliaSannino2024}:
\begin{equation}\label{eq:scalar-mass-scale}
  \mu^2 \;\sim\; \xi\, \frac{\Lambda_D^4}{M_{\text{Pl}}^2}\,,
\end{equation}
(Eq.~(22) of~\cite{CacciapagliaSannino2024}), where $\xi$ is an
$\mathcal{O}(1)$ numerical coefficient depending on the $c^{abcd}$
couplings and on form factors connecting the four-fermion operators to
the composite scalar masses.

\subsection{Induced VEVs for the dormant Higgs doublets}

The full scalar potential for the 18 Higgs doublets takes the
form~\cite{CacciapagliaSannino2024}:
\begin{equation}\label{eq:full-potential}
  V_H \;=\; \lambda\, (H_0^\dagger H_0)^2
  \;-\; \mu_0^2\, (H_0^\dagger H_0)
  \;-\; \mu_{0i}^2\, (H_0^\dagger H_i + H_i^\dagger H_0)
  \;+\; M_{ij}^2\, (H_i^\dagger H_j)\,,
\end{equation}
(Eq.~(17) of~\cite{CacciapagliaSannino2024}), where $i,j = 1,\ldots,17$
label the dormant Higgs doublets and the quartic self-couplings of the
dormant doublets have been neglected (their VEVs are suppressed).

The key point is that the leading Higgs $H_0$ gets a
\textbf{negative mass-squared} $-\mu_0^2 < 0$ (from radiative
corrections or from the strong dynamics at the duality scale) and
triggers electroweak symmetry breaking.
The selection of $H_0$ as the direction with negative mass-squared is a dynamical input to the framework (from radiative corrections or strong-sector dynamics), not a result derived within this lecture.  The dormant Higgs doublets
have large \textbf{positive mass-squared} $M_{ij}^2 > 0$ and do not
break the electroweak symmetry on their own.  However, the
\textbf{off-diagonal mixing terms} $\mu_{0i}^2$ induce small VEVs for
the dormant doublets:
\begin{equation}\label{eq:induced-vev}
  \boxed{%
  \langle H_i \rangle
  \;=\; (M_{ij}^2)^{-1}\, \mu_{0i}^2\, \langle H_0 \rangle\,,}
\end{equation}
(Eq.~(18) of~\cite{CacciapagliaSannino2024}).
This is simply the result of integrating out the 17 heavy Higgs fields
at tree level.  The induced VEVs are suppressed relative to
$\langle H_0 \rangle$ by the ratio $\mu_{0i}^2 / M_{ij}^2$, which is
controlled by the unknown $\mathcal{O}(1)$ coefficients in the
Planck-scale operators.

\subsection{VEV tiers and fermion mass hierarchy}

Following~\cite{CacciapagliaSannino2024}, the VEV hierarchy can be
organised into ``tiers.''  The first tier is the leading VEV
$\langle H_0 \rangle \approx v$, which gives the top quark its mass.
The second tier arises from mixing between $H_0$ and a set of dormant
Higgs doublets with mass $M_{d,33}$ (those coupling to the third
generation down-type quarks).

The mixing is generated by Planck-scale operators of the
form~\cite{CacciapagliaSannino2024}:
\begin{equation}\label{eq:mixing-2nd-tier}
  -\mu_b^2\, H_{33}^{u\,\dagger} H_{33}^d
  \;-\; \mu_c^2\, H_{33}^{u\,\dagger} (H_{22}^u)^\dagger
  \;-\; \mu_{23}^2\, H_{33}^{u\,\dagger} (H_{23}^u + H_{32}^u)^\dagger
  \;+\; \text{h.c.}
\end{equation}
(Eq.~(13) of~\cite{CacciapagliaSannino2024}).  The resulting VEVs at
second tier are:
\begin{equation}\label{eq:2nd-tier-vev}
  \text{(2nd tier)}\qquad
  \langle H_{33}^d \rangle
  \;\sim\; \frac{\mu^2}{M_{d,33}^2}\, v
  \;\sim\; \frac{m_b}{m_t}\, v\,,
\end{equation}
(Eq.~(14) of~\cite{CacciapagliaSannino2024}),
and similarly for the charm and strange quark masses, which arise from
mixing between the third and second generation Higgs doublets.  A third
tier produces even smaller VEVs:
\begin{equation}\label{eq:3rd-tier-vev}
  \text{(3rd tier)}\qquad
  \langle H \rangle
  \;\sim\; \frac{\mu^4}{M^4}\, v\,,
\end{equation}
(Eq.~(15) of~\cite{CacciapagliaSannino2024}),
and so on, generating a cascade of progressively smaller VEVs.

The pattern is clear: \emph{if the duality holds}, the fermion mass
hierarchy is generated by a hierarchy of induced VEVs, each tier
further suppressed by a factor of $\mu^2/M^2 \ll 1$.  The
$\mathcal{O}(1)$ coefficients $c^{abcd}$ in the Planck-scale operators
determine the precise pattern of suppression, but the
\emph{parametric structure}---geometrically decreasing VEVs, one tier
per power of $\mu^2/M^2$---is a consequence of the framework.

\begin{warningbox}[title={Free parameters in the mass hierarchy}]
  The coefficients $c^{abcd}$ in Eq.~\eqref{eq:planck-operator} are
  \textbf{free $\mathcal{O}(1)$ parameters}.  The dualised SM does
  \emph{not} predict the specific values of fermion masses---it provides
  a \textbf{framework} where the hierarchy is \emph{natural} (arising
  from a VEV hierarchy driven by Planck-scale operators) rather than
  \emph{put in by hand} (via a hierarchy in Yukawa couplings).

  \medskip

  The framework has $\mathcal{O}(100)$ free parameters (the $c^{abcd}$
  tensor), comparable to the number of free parameters in the SM
  Yukawa sector (quark and lepton masses, mixing angles, and CP phases).
  The conceptual gain is that these parameters have a
  \emph{physical origin}---Planck-scale gravity-induced
  operators---rather than being arbitrary Lagrangian couplings.
\end{warningbox}

\subsection{Lepton masses}

Lepton masses can be generated via the same mechanism, \emph{if} the
leptons participate in the duality via the Pati--Salam extension
(lepton as fourth colour, discussed in Lecture~7, \S2).  In the QCD-sector duality where leptons are
elementary, an additional four-fermion operator is
required~\cite{CacciapagliaSannino2024}:
\begin{equation}\label{eq:lepton-operator}
  \mathcal{L}_\ell
  \;\supset\;
  \frac{h_l}{M^2}\, (L\widetilde{L})(Q\widetilde{Q})^\dagger\,,
\end{equation}
(Eq.~(24) of~\cite{CacciapagliaSannino2024}),
which at the duality scale generates a direct coupling of the leptons
to the mes-Higgs field:
\begin{equation}\label{eq:lepton-yukawa}
  \mathcal{L}_m
  \;\supset\;
  h_l\, \mathcal{F}_\Phi\,
  \frac{\Lambda_D^2}{M^2}\,
  \widetilde{l}\, l\, \Phi_H^\dagger\,,
\end{equation}
(Eq.~(25) of~\cite{CacciapagliaSannino2024}),
where $\mathcal{F}_\Phi$ is a form factor.  The expansion in
terms of the individual Higgs doublets gives:
\begin{equation}\label{eq:lepton-yukawa-expanded}
  \widetilde{l}\, l\, \Phi_H^\dagger
  \;=\; \sum_{i,j}
  \Bigl(
    \ell_L^i\, \nu_R^j\, (H_{ij}^d)^\dagger
    \;+\; \ell_L^i\, e_R^j\, (H_{ij}^u)^\dagger
  \Bigr)\,,
\end{equation}
(Eq.~(26) of~\cite{CacciapagliaSannino2024}).
Interestingly, the largest VEV couples naturally to the $\tau$ lepton,
explaining its relatively large mass compared to the electron and muon.
The hierarchy between charged lepton masses follows the same VEV tier
pattern as the quark sector.

Large masses for right-handed neutrinos can be generated close to the
duality scale $\Lambda_D$, implementing a type-I seesaw
mechanism~\cite{CacciapagliaSannino2024}.


%% ========================================================================
%%  SECTION 3: Electroweak Symmetry Breaking in the Multi-Higgs Sector
%% ========================================================================
\section{Electroweak Symmetry Breaking in the Multi-Higgs Sector}
\label{sec:ewsb}

The presence of 18 Higgs doublets in the magnetic theory raises an
immediate question: how does electroweak symmetry breaking (EWSB) occur
in such a rich scalar sector, and how is the observed SM-like Higgs
boson (the 125~GeV scalar discovered at the LHC) recovered at low
energies?

\subsection{One light Higgs, seventeen heavy}

The 18 Higgs doublets have a scalar potential determined by the
gauge and global symmetries of the magnetic theory, plus the
Planck-scale operators of Section~\ref{sec:planck-operators}.  The
essential structure is:
\begin{itemize}
  \item \textbf{One linear combination}
    ($H_0 \equiv H_{33}^u$) gets a \textbf{negative mass-squared}
    $-\mu_0^2$ from radiative corrections (Coleman--Weinberg mechanism)
    or from the strong dynamics at $\Lambda_D$, and triggers EWSB:
    $\langle H_0 \rangle \approx v$.
  \item The remaining \textbf{17 dormant Higgs doublets}
    have \textbf{positive mass-squared} of order
    $M_i^2 \sim \Lambda_D^2 / M_{\text{Pl}}$ or larger, keeping them
    heavy.  Their induced VEVs (Eq.~\eqref{eq:induced-vev}) are
    hierarchically smaller than $v$.
\end{itemize}

The effective single-Higgs theory that emerges below the mass scale of
the heavy doublets reproduces the SM Higgs potential:
\begin{equation}\label{eq:sm-potential}
  V_{\text{eff}} \;=\;
  \lambda\, (H_0^\dagger H_0)^2
  \;-\; m^2\, (H_0^\dagger H_0)\,,
\end{equation}
(Eq.~(19) of~\cite{CacciapagliaSannino2024}), with
\begin{equation}\label{eq:effective-mass}
  m^2 \;=\; \mu_0^2
  \;+\; \sum_{i,j=1}^{17}
  \mu_{0j}^2\, (M_{ij}^2)^{-1}\, \mu_{0i}^2\,,
\end{equation}
(Eq.~(20) of~\cite{CacciapagliaSannino2024}).
The first term is the ``bare'' Higgs mass-squared, and the second is
the correction from integrating out the 17 heavy doublets (a
positive-definite shift, since the mixing masses $\mu_{0i}^2$
contribute constructively).

\subsection{The VEV sum rule as a constraint}

The total electroweak VEV is distributed among all 18 doublets:
\begin{equation}\label{eq:vev-constraint}
  v^2 \;=\; \sum_{i,j=1}^{3}
  \Bigl(
    \langle H_{ij}^u \rangle^2
    \;+\; \langle H_{ij}^d \rangle^2
  \Bigr)
  \;=\; \frac{1}{2}\,(246\;\text{GeV})^2\,.
\end{equation}
Since the leading VEV dominates ($\langle H_0 \rangle \gg \langle H_i
\rangle$ for $i \neq 0$), the VEV sum rule is satisfied to a good
approximation by the leading term alone:
\begin{equation}\label{eq:vev-dominance}
  \langle H_0 \rangle^2 \;\approx\; v^2
  \;-\; \sum_{i \neq 0} \langle H_i \rangle^2
  \;\approx\; v^2
  \Bigl(1 - \mathcal{O}\bigl(m_b^2/m_t^2\bigr)\Bigr)
  \;\approx\; v^2\,,
\end{equation}
where we used the VEV hierarchy
$\langle H_i \rangle / v \sim m_b/m_t \sim 0.024$ for the
largest dormant VEV.  The corrections are at the per-mille level.

\begin{remark}[The VEV sum rule is an input]
  We emphasise again that the VEV sum rule~\eqref{eq:vev-constraint}
  is \emph{not} a prediction of the dualised SM.  It is a consequence
  of the measured Fermi constant $G_F$, which determines the
  electroweak scale $v = (\sqrt{2}\, G_F)^{-1/2} = 246\;\text{GeV}$.
  The dualised SM provides a mechanism for \emph{distributing} this
  measured VEV among the 18 doublets, but the total is an input.
\end{remark}

\subsection{Low-energy effective theory: one Higgs doublet}

At energies $E \ll M_i$ (where $M_i$ are the masses of the 17 heavy
doublets), the magnetic theory reduces to an effective theory with
\textbf{one light Higgs doublet} $H_0$ plus the SM fermions and gauge
bosons.  This is precisely the Standard Model.

The 125~GeV scalar discovered at the LHC is identified with the
physical Higgs boson---the neutral CP-even component of $H_0$ after
EWSB.  \emph{If the duality holds}, this scalar is \emph{not}
elementary: it is a composite state
$\Phi_H = \langle \widetilde{Q}\lambda\lambda Q \rangle$ of the
electric theory (Lecture~7, Eq.~(3) of~\cite{CacciapagliaSannino2024}).
The compositeness becomes manifest only at the
duality scale $\Lambda_D \sim 10^{11}\;\text{GeV}$, far above current
collider energies.

The remaining 17 Higgs doublets, 17 mesinos (the fermionic partners of
the dormant Higgs fields), and the additional quasi-SUSY states
(gaugino, squarks) are all heavy, with masses at or above the scale
$\Lambda_S \sim \text{few TeV}$.  Their effects on low-energy
observables are suppressed by powers of $v / \Lambda_S$ and are
consistent with current precision electroweak measurements.

\begin{previewbox}[title={Preview: Exercises 5 and 6}]
  Exercises~5 and 6 (distributed separately) make the parametric
  statements from this and the following sections quantitative:
  \begin{itemize}
    \item \textbf{Exercise~5:} Given the universal Yukawa coupling $y$,
      the VEV hierarchy from Planck-scale operators, and specific
      $\mathcal{O}(1)$ coefficients, compute the fermion mass ratios
      and show parametric consistency with observations.
    \item \textbf{Exercise~6:} Given the duality scale
      $\Lambda_D = 10^{11}\;\text{GeV}$, extract
      $\Lambda_{\QCD} \approx 200\;\text{MeV}$ from one-loop RG running
      and verify $v \sim \Lambda_D^2 / M_{\text{Pl}} \sim
      \mathcal{O}(\text{TeV})$.
  \end{itemize}
\end{previewbox}


%% ========================================================================
%%  SECTION 4: Scale Matching and Physical Parameters
%% ========================================================================
\section{Scale Matching and Physical Parameters}
\label{sec:scale-matching}

The dualised SM framework, \emph{if the duality holds}, connects the
high-energy dynamics of the electric theory to the measured parameters
of the SM at low energies.  In this section, we derive the parametric
relationships between the duality scale $\Lambda_D$, the Fermi scale
$v$, and the QCD scale $\Lambda_{\QCD}$.

\subsection{The duality scale}

The duality scale $\Lambda_D$ is the energy at which the electric gauge
coupling becomes strong and the electric theory confines.  Below
$\Lambda_D$, the magnetic (SM) description is appropriate.

From Eq.~\eqref{eq:scalar-mass-scale}, the scalar mass-squared is
$\mu^2 \sim \xi \,\Lambda_D^4 / M_{\text{Pl}}^2$.  Requiring that the
scalar mass lies near the electroweak scale,
$\mu \sim \mathcal{O}(\text{TeV})$, gives an estimate of the duality
scale~\cite{CacciapagliaSannino2024}:
\begin{equation}\label{eq:duality-scale}
  \boxed{%
  \xi^{-1/4}\, \Lambda_D
  \;\sim\; \sqrt{\mu\, M_{\text{Pl}}}
  \;\sim\; 10^{11}\;\text{GeV}\,,}
\end{equation}
(Eq.~(23) of~\cite{CacciapagliaSannino2024}), where we used
$\mu \sim 1\;\text{TeV}$ and
$M_{\text{Pl}} \sim 10^{19}\;\text{GeV}$.

This is a \emph{seesaw-like} relation: the geometric mean of the TeV
scale and the Planck scale gives an intermediate scale of
$10^{11}\;\text{GeV}$.  The $\mathcal{O}(1)$ coefficient $\xi$ (from
the Planck-scale operators) introduces an uncertainty of perhaps one
order of magnitude in either direction.

\begin{remark}[$\Lambda_D$ is a free parameter]
  The duality scale $\Lambda_D$ is \emph{not} predicted by the
  framework.  It is a free parameter---analogous to
  $\Lambda_{\QCD}$ in ordinary QCD---that must ultimately be determined
  from experiment.  The estimate $\Lambda_D \sim 10^{11}\;\text{GeV}$
  is a parametric consistency check: given the measured electroweak
  scale and the Planck scale, this is the \emph{natural} value of
  $\Lambda_D$ within the framework.
\end{remark}

\subsection{The Fermi scale from the duality scale}

Inverting the seesaw relation~\eqref{eq:duality-scale}, we can express
the electroweak VEV $v$ in terms of $\Lambda_D$ and $M_{\text{Pl}}$:
\begin{equation}\label{eq:fermi-scale}
  v^2 \;\sim\; \mu^2
  \;\sim\; \frac{\Lambda_D^4}{M_{\text{Pl}}^2}\,,
  \qquad
  v \;\sim\; \frac{\Lambda_D^2}{M_{\text{Pl}}}
  \;\sim\; \frac{(10^{11}\;\text{GeV})^2}{10^{19}\;\text{GeV}}
  \;=\; 10^3\;\text{GeV}\,,
\end{equation}
which is parametrically consistent with the measured value
$v = 246\;\text{GeV}$.  The factor of $\sim 4$ discrepancy between
$10^3\;\text{GeV}$ and $246\;\text{GeV}$ is absorbed by the
$\mathcal{O}(1)$ coefficient $\xi$.

\begin{warningbox}[title={Parametric consistency, not prediction}]
  The relation $v \sim \Lambda_D^2 / M_{\text{Pl}}$ is a
  \textbf{dimensional analysis argument}.  It shows that the
  \emph{parametric dependence} of $v$ on $\Lambda_D$ and
  $M_{\text{Pl}}$ is correct---the electroweak hierarchy
  $v / M_{\text{Pl}} \sim 10^{-17}$ follows from the intermediate
  scale $\Lambda_D / M_{\text{Pl}} \sim 10^{-8}$ via a seesaw
  mechanism.  The \emph{numerical value} $v = 246\;\text{GeV}$ is not
  predicted; it depends on the $\mathcal{O}(1)$ coefficient $\xi$,
  which is a free parameter.
\end{warningbox}

\subsection{$\Lambda_{\QCD}$ from one-loop RG running}
\label{ssec:lambda-qcd}

A striking consistency check of the framework is that the QCD
confinement scale $\Lambda_{\QCD} \approx 200\;\text{MeV}$ can be
extracted from the one-loop running of the strong coupling constant
$\alphas$, starting from its value at the duality scale $\Lambda_D$.

At the duality scale, the magnetic and electric gauge couplings are
related by~\cite{CacciapagliaSannino2024}:
\begin{equation}\label{eq:coupling-duality}
  \frac{1}{\alpha_e(\Lambda_D)}
  \;\sim\; \alpha_m(\Lambda_D)
  \;=\; \frac{g_m^2}{4\pi}\,,
\end{equation}
(Eq.~(9) of~\cite{CacciapagliaSannino2024}).
Strong coupling in the electric theory ($\alpha_e \gg 1$) corresponds
to weak coupling in the magnetic theory ($\alpha_m \ll 1$).

The one-loop beta function for the magnetic $\SU{3}_c$ gauge coupling
gives the standard relation between the UV scale and the confinement
scale (from Lecture~1):
\begin{equation}\label{eq:rg-running}
  \Lambda_{\QCD}
  \;=\; \mu_{\UV} \;\exp\!\Bigl(
    -\frac{2\pi}{\bzero\,\alphas(\mu_{\UV})}
  \Bigr)\,,
\end{equation}
where $\bzero = (11\Nc - 2\Nf)/3$ is the one-loop beta function
coefficient.  The crucial observation is that $\Lambda_{\QCD}$ is
\textbf{exponentially sensitive} to $\alphas(\mu_{\UV})$ and to
$\bzero$: small changes in the input coupling produce large changes in
$\Lambda_{\QCD}$.

To illustrate the parametric consistency, we work \emph{backwards} from
the known SM.  The measured value
$\alphas(M_Z) = 0.118$ at $M_Z = 91.2\;\text{GeV}$ can be evolved to
$\Lambda_D = 10^{11}\;\text{GeV}$ using the one-loop formula
with $\bzero = 7$ ($\Nf = 6$):
\begin{equation}\label{eq:alpha-running}
  \alphas(\Lambda_D)
  \;=\; \frac{\alphas(M_Z)}{1 + \dfrac{\bzero}{2\pi}\,
    \alphas(M_Z)\, \ln\!\bigl(\Lambda_D / M_Z\bigr)}
  \;\approx\; \frac{0.118}{1 + \dfrac{7}{2\pi} \times 0.118 \times 20.8}
  \;\approx\; 0.032\,.
\end{equation}
Substituting back into Eq.~\eqref{eq:rg-running} with
$\mu_{\UV} = \Lambda_D$:
\begin{equation}\label{eq:lambda-qcd-estimate}
  \Lambda_{\QCD}
  \;\sim\; 10^{11}\;\text{GeV} \times
  \exp\!\Bigl(-\frac{2\pi}{7 \times 0.032}\Bigr)
  \;=\; 10^{11}\;\text{GeV} \times
  \exp(-28.0)
  \;\sim\; 10^{11}\;\text{GeV} \times
  6 \times 10^{-13}\,,
\end{equation}
giving
\begin{equation}\label{eq:lambda-qcd-result}
  \boxed{\Lambda_{\QCD}^{(\Nf=6)}
  \;\sim\; 60\;\text{MeV}\,.}
\end{equation}
This is within a factor of 3 of the physical value
$\Lambda_{\QCD} \approx 200\;\text{MeV}$---the discrepancy is expected,
because (i)~the number of active flavours decreases at each quark mass
threshold (from $\Nf = 6$ above $m_t$ down to $\Nf = 3$ below $m_c$),
which increases $\bzero$ and slows the running, and (ii)~the precise
matching of couplings at $\Lambda_D$ depends on the $\mathcal{O}(1)$
duality relation~\eqref{eq:coupling-duality}, which is not
quantitatively determined.

The point of this exercise is \textbf{not} to derive
$\Lambda_{\QCD} = 200\;\text{MeV}$ from first principles.  It is to
show that the one-loop QCD running, starting from a coupling
$\alphas(\Lambda_D) \sim 0.03$ at
$\Lambda_D \sim 10^{11}\;\text{GeV}$, produces a confinement scale in
the \emph{right parametric window}---between $\sim 50\;\text{MeV}$ and
$\sim 1\;\text{GeV}$, depending on the input coupling and threshold
treatment.

\begin{remark}[This is not a precision prediction]
  The estimate $\Lambda_{\QCD} \sim 200\;\text{MeV}$ depends on the
  input coupling $\alphas(\Lambda_D)$, which in the dualised SM
  framework is determined by the duality relation at $\Lambda_D$.  The
  value $\alphas(\Lambda_D) \sim 0.03$ is obtained by running the
  measured $\alphas(M_Z)$ up to $10^{11}\;\text{GeV}$.  This is a
  \textbf{consistency check}: the framework is compatible with the
  measured QCD scale if $\Lambda_D \sim 10^{11}\;\text{GeV}$ and the
  duality coupling is in a reasonable range.  A true prediction would
  require calculating $\alphas(\Lambda_D)$ from the duality relation
  without using the measured $\alphas(M_Z)$ as input.
\end{remark}


%% ========================================================================
%%  SECTION 5: Three Scales and Collider Signatures
%% ========================================================================
\section{Three Scales and Collider Signatures}
\label{sec:collider}

\subsection{The energy hierarchy}

\emph{If the duality holds}, the physics of the dualised SM is
characterised by a hierarchy of energy scales (recalled from Lecture~7,
\S4.3):
\begin{enumerate}
  \item $\Lambda_S \sim \text{few TeV}$: the \textbf{quasi-SUSY
    spectrum} appears---gaugino, squarks, heavy Higgs doublets, and
    mesino states with masses at or above $\Lambda_S$.  Below this
    scale, only SM degrees of freedom are active.
  \item $\Lambda_D \sim 10^{11}\;\text{GeV}$: the \textbf{duality
    scale}---the energy at which the electric $\SU{2}$ gauge coupling
    becomes strong and the electric theory confines.  The magnetic SM
    description is valid below $\Lambda_D$.
  \item $\Lambda_{\text{eGUT}}$ (optional): if the three gauge
    couplings of the electric theory unify at a scale
    $\Lambda_{\text{eGUT}} < M_{\text{Pl}}$, a \textbf{grand unified
    theory of the electric sector} emerges.  This is an attractive but
    speculative possibility.
\end{enumerate}
The running of the gauge couplings through these three regimes is
illustrated schematically in Fig.~1 of~\cite{CacciapagliaSannino2024}:
the SM couplings run normally below $\Lambda_S$; between $\Lambda_S$
and $\Lambda_D$, the magnetic beta functions are modified by the
quasi-SUSY spectrum; and above $\Lambda_D$, the electric couplings run
with only fermion and gauge boson contributions (no scalars).

\subsection{Collider signatures at $\Lambda_S$}

The quasi-SUSY spectrum of the dualised SM predicts new states at the
scale $\Lambda_S \sim \text{few TeV}$~\cite{CacciapagliaSannino2024}.
The most accessible at hadron colliders are:

\paragraph{Gaugino-like state.}
The magnetic gaugino $\lambda_m$---a colour-octet Weyl
fermion---behaves like the \textbf{gluino} of the Minimal
Supersymmetric Standard Model (MSSM) from a collider perspective.  It
is pair-produced via strong interactions ($gg \to \lambda_m
\overline{\lambda}_m$), and its decay products include jets and
missing transverse energy (if the lightest magnetic state is stable on
detector timescales).  Current ATLAS and CMS searches constrain the
gluino mass to be above roughly
$2\;\text{TeV}$~\cite{CacciapagliaSannino2024}, and these bounds
apply with caveats to the quasi-SUSY gaugino as well.

\paragraph{Squark-like states.}
The scalar colour triplets $\phi$ and $\widetilde{\phi}$---the
quasi-SUSY partners of the magnetic quarks---behave like
\textbf{squarks} in collider searches.  They are pair-produced via
strong interactions and decay to quarks plus the gaugino (if
kinematically allowed) or via gauge-mediated channels.  The mass limits
from LHC searches are comparable to those for MSSM squarks.

\paragraph{Heavy Higgs doublets and mesino.}
The 17 dormant Higgs doublets and the mesino $M$ are electroweak states
(colour singlets or doublets under $\SU{2}_L$).  Their production
cross-sections at hadron colliders are smaller than those of the
strongly interacting states, but they could be accessible at future
high-energy colliders or through precision flavour
observables~\cite{CacciapagliaSannino2024}.

If the lightest mesino component has a Majorana mass, it could serve as
a dark matter candidate~\cite{CacciapagliaSannino2024}.

\begin{warningbox}[title={Key discriminator: quasi-SUSY vs.\ true SUSY}]
  If a spectrum resembling the MSSM is discovered at the LHC or a
  future collider, how would we distinguish the dualised SM from true
  supersymmetry?  The key differences are:
  \begin{enumerate}
    \item \textbf{The Higgs is composite}, not elementary.  In the
      MSSM, the Higgs doublets $H_u$ and $H_d$ are elementary chiral
      superfields.  In the dualised SM, the Higgs is a composite meson
      $\Phi_H = \langle \widetilde{Q}\lambda\lambda Q\rangle$, with
      compositeness effects becoming visible at $\Lambda_D$.
    \item \textbf{18 Higgs doublets}, not 2.  The MSSM has two Higgs
      doublets; the dualised SM has 18.  The 17 heavy doublets have
      masses $\sim \Lambda_S$--$\Lambda_D$.
    \item \textbf{No superpartner of the Higgs}.  In the MSSM, the
      higgsino is a fermionic partner of the Higgs.  In the dualised
      SM, the fermionic partner of the mes-Higgs is the \emph{mesino}
      $M$, which transforms differently under the flavour symmetry.
    \item \textbf{No $R$-parity from fundamental symmetry.}  In the
      MSSM, $R$-parity distinguishes SM particles from superpartners.
      In the dualised SM, the quasi-SUSY pairing is emergent, not
      fundamental, and any stabilising symmetry for the lightest new
      state must be argued separately.
  \end{enumerate}
\end{warningbox}


\subsection{Causal chain summary}

The following table collects the logical chain developed in this lecture,
from the Planck-scale input to the measured low-energy parameters.

\begin{center}
\renewcommand{\arraystretch}{1.3}
\begin{tabular}{@{}l l l@{}}
\toprule
\textbf{Step} & \textbf{What it determines} & \textbf{Free parameters} \\
\midrule
Planck operators (\S\ref{sec:planck-operators})
  & Four-fermion couplings $\to$ scalar masses
  & $c^{abcd}$ ($\mathcal{O}(1)$ coefficients) \\
Scalar mixing (\S\ref{sec:ewsb})
  & Induced VEVs of dormant Higgs doublets
  & $\mu_{0i}^2 / M_i^2$ ratios \\
Universal Yukawa (\S\ref{sec:scalar-democracy})
  & Fermion mass hierarchy (tier structure)
  & VEV pattern from above \\
RG running (\S\ref{sec:scale-matching})
  & $v$, $\Lambda_{\QCD}$ from $\Lambda_D$
  & $\Lambda_D$ (single input scale) \\
\bottomrule
\end{tabular}
\end{center}


%% ========================================================================
%%  SECTION 6: Summary and Honest Assessment
%% ========================================================================
\section{Summary and Honest Assessment}
\label{sec:summary-assessment}

This lecture has developed the quantitative consequences of the
dualised SM framework introduced in Lecture~7.  We conclude with an
honest assessment of what the programme achieves and what remains open.

\subsection{What the dualised SM does}

\emph{If the conjectured non-SUSY duality holds}, the dualised SM
provides:
\begin{enumerate}
  \item \textbf{An alternative UV description of the SM with no
    elementary scalars.}  The electric theory
    ($\SU{2}$ with 6~Dirac flavours and one adjoint Weyl fermion)
    contains only fermions and gauge bosons.  All scalars in the
    low-energy (magnetic) SM description are composites
    (Lecture~7, \S1.3).

  \item \textbf{A natural resolution of the hierarchy problem.}
    Since there are no elementary scalars in the electric theory,
    there is no scalar mass to protect against quadratic divergences.
    The Higgs mass is set by the confinement dynamics at $\Lambda_D$,
    just as the pion mass in QCD is set by chiral symmetry breaking,
    not by a fundamental scalar mass parameter (Lecture~7, \S4.1).

  \item \textbf{A framework for the fermion mass hierarchy.}
    The scalar democracy mechanism
    (Section~\ref{sec:scalar-democracy})---universal Yukawa coupling
    $y$ combined with hierarchical Higgs VEVs induced by Planck-scale
    operators (Section~\ref{sec:planck-operators})---generates the
    observed pattern of fermion masses without requiring a hierarchy in
    the Yukawa couplings.  The $\mathcal{O}(1)$ coefficients $c^{abcd}$
    are free parameters, but their \emph{origin} (Planck-scale gravity)
    is specified.

  \item \textbf{Parametric consistency with the electroweak and QCD
    scales.}  The duality scale $\Lambda_D \sim 10^{11}\;\text{GeV}$
    gives
    $v \sim \Lambda_D^2/M_{\text{Pl}} \sim \mathcal{O}(\text{TeV})$
    (Section~\ref{sec:scale-matching}) and
    $\Lambda_{\QCD} \sim 200\;\text{MeV}$ from one-loop RG running
    (\S\ref{ssec:lambda-qcd}).

  \item \textbf{A testable quasi-SUSY spectrum at the TeV scale.}
    The gaugino, squarks, and heavy Higgs doublets predicted by the
    magnetic theory (Section~\ref{sec:collider}) are accessible to
    current and future collider experiments.

  \item \textbf{Three or more generations as a structural
    requirement.}  The generation bound $n_g \geq 3$
    (Lecture~7, \S2.4) follows from requiring the electric gauge group
    to be non-abelian.
\end{enumerate}

\subsection{What the dualised SM does not do}

\begin{enumerate}
  \item \textbf{It does not prove the non-SUSY duality.}
    The duality is a \emph{conjecture}, supported by anomaly matching
    (Lecture~7, \S5) but without a rigorous proof analogous to
    Seiberg's SUSY result (Lecture~3).  The central vulnerability,
    articulated in Lecture~6, remains: anomaly matching is necessary but
    not sufficient.

  \item \textbf{It does not predict specific fermion masses.}
    The $\mathcal{O}(1)$ coefficients $c^{abcd}$ in the Planck-scale
    operators are free parameters.  The framework explains why fermion
    masses are \emph{hierarchical} but does not predict their
    \emph{values}.

  \item \textbf{It does not predict the duality scale $\Lambda_D$.}
    The estimate $\Lambda_D \sim 10^{11}\;\text{GeV}$ is a consistency
    argument, not a first-principles calculation.  $\Lambda_D$ is a
    free parameter analogous to $\Lambda_{\QCD}$ in QCD.

  \item \textbf{It does not resolve the Cheng--Shadmi tension.}
    As discussed in Lecture~6 (\S5), the fate of the duality under
    large SUSY breaking (required for the non-SUSY version) remains an
    open theoretical question.  The Cheng--Shadmi counterexamples
    caution against assuming that SUSY duality survives arbitrary
    deformations.
\end{enumerate}

\subsection{The pedagogical arc: Lectures 1--8}

These eight lectures have traced the following path:
\begin{itemize}
  \item \textbf{Lectures 1--2:} Foundations of anomaly matching and
    supersymmetric gauge theories.  Anomaly matching as the central
    tool; superfields and the SQCD Lagrangian.
  \item \textbf{Lectures 3--4:} Seiberg duality in $\mathcal{N} = 1$
    SQCD.  Electric--magnetic duality with explicit anomaly
    verification (7 out of 7 conditions), the conformal window, mass
    deformations, and the connection to $\mathcal{N} = 2$ theories
    via the Argyres--Plesser--Seiberg deformation.
  \item \textbf{Lectures 5--6:} Non-perturbative dynamics beyond SUSY.
    Tumbling and the MAC criterion, complementarity, and the
    non-SUSY duality programme---its history, evidence, and central
    vulnerability.
  \item \textbf{Lectures 7--8:} The dualised Standard Model.  The
    complete construction (electric theory, magnetic spectrum, anomaly
    matching) and its quantitative consequences (scalar democracy,
    scale matching, collider signatures).
\end{itemize}

\subsection{Testability and outlook}

The programme is testable: the quasi-SUSY spectrum at the TeV scale is
a concrete prediction that LHC and future collider experiments can
confirm or exclude.  Specifically:
\begin{itemize}
  \item If \textbf{no quasi-SUSY states} are found below
    $\sim 10\;\text{TeV}$, the simplest version of the dualised SM
    (with $\Lambda_S \sim \text{few TeV}$) is excluded, though
    variants with higher $\Lambda_S$ remain viable.
  \item If \textbf{quasi-SUSY states are found}, the key question is
    whether the Higgs is elementary (as in the MSSM) or composite (as
    in the dualised SM).  Precision measurements of Higgs couplings,
    deviations from SM predictions, and the spectrum of heavy Higgs
    doublets could distinguish between the two scenarios.
  \item Lattice simulations of the electric theory
    ($\SU{2}$ with 6~fundamentals and one adjoint) could in principle
    test the duality non-perturbatively, though this is a technically
    challenging computation.
\end{itemize}

\begin{tcolorbox}[
  enhanced,
  colback=green!5!white,
  colframe=green!60!black,
  fonttitle=\bfseries,
  title={Lecture 8: Summary}
]
\begin{enumerate}
  \item \textbf{Scalar democracy:}
    A universal magnetic Yukawa coupling $y$ combined with hierarchical
    Higgs VEVs generates the fermion mass hierarchy.  The flavour
    puzzle is translated from the Yukawa sector to the scalar (VEV)
    sector~(\S\ref{sec:scalar-democracy}).
  \item \textbf{Planck-scale operators:}
    Four-fermion operators of the form
    $c^{abcd}(Q^a\widetilde{Q}^b)(Q^c\widetilde{Q}^d)^\dagger
    / M_{\text{Pl}}^2$ generate the VEV hierarchy.  The $\mathcal{O}(1)$
    coefficients $c^{abcd}$ are free
    parameters~(\S\ref{sec:planck-operators}).
  \item \textbf{EWSB:}
    One Higgs doublet ($H_0 = H_{33}^u$) triggers EWSB; the remaining
    17 are heavy with induced VEVs.  The low-energy effective theory is
    the SM with one composite Higgs~(\S\ref{sec:ewsb}).
  \item \textbf{Scale matching:}
    $\Lambda_D \sim 10^{11}\;\text{GeV}$ gives
    $v \sim \Lambda_D^2/M_{\text{Pl}} \sim \mathcal{O}(\text{TeV})$ and
    $\Lambda_{\QCD} \sim 200\;\text{MeV}$ from one-loop RG running.
    These are consistency checks, not
    predictions~(\S\ref{sec:scale-matching}).
  \item \textbf{Testability:}
    Quasi-SUSY states at $\Lambda_S \sim \text{TeV}$ are the smoking
    gun.  Key discriminators from true SUSY: composite Higgs, 18
    doublets, no fundamental
    $R$-parity~(\S\ref{sec:collider}).
  \item \textbf{Honest assessment:}
    The framework provides parametric consistency and a testable
    spectrum.  The duality itself remains conjectural, and the
    $\mathcal{O}(1)$ coefficients are
    free~(\S\ref{sec:summary-assessment}).
\end{enumerate}
\end{tcolorbox}

\begin{previewbox}[title={Exercises 5 and 6}]
  Exercises~5 and 6 make the parametric statements from this lecture
  quantitative.  Exercise~5 works through the VEV hierarchy for the
  fermion mass spectrum; Exercise~6 extracts
  $v \sim 246\;\text{GeV}$ and
  $\Lambda_{\QCD} \sim 200\;\text{MeV}$ from $\Lambda_D$ using
  one-loop RG running.  These exercises are distributed separately.
\end{previewbox}


%% ========================================================================
%%  BIBLIOGRAPHY
%% ========================================================================
\begin{thebibliography}{99}

\bibitem{CacciapagliaSannino2024}
  G.~Cacciapaglia and F.~Sannino,
  ``Charting Standard Model duality and its signatures,''
  \textit{Phys.~Rev.} \textbf{D111} (2025) 035013
  [\texttt{arXiv:2407.17281}].

\bibitem{HillMachado2019}
  C.~T.~Hill, P.~A.~N.~Machado, A.~E.~Thomsen, and J.~Turner,
  ``Scalar democracy,''
  \textit{Phys.~Rev.} \textbf{D100} (2019) 015015
  [\texttt{arXiv:1902.07214}].

\bibitem{Sannino2011}
  F.~Sannino,
  ``The Standard Model is natural as magnetic gauge theory,''
  \textit{Mod.~Phys.~Lett.} \textbf{A26} (2011) 1763
  [\texttt{arXiv:1102.5100}].

\bibitem{PS1974}
  J.~C.~Pati and A.~Salam,
  ``Lepton number as the fourth `color',''
  \textit{Phys.~Rev.} \textbf{D10} (1974) 275.

\bibitem{AMPS2011}
  O.~Antipin, M.~Mojaza, C.~Pica, and F.~Sannino,
  ``Magnetic fixed points and emergent supersymmetry,''
  \textit{JHEP} \textbf{1306} (2013) 037
  [\texttt{arXiv:1105.1510}].

\bibitem{Seiberg1995}
  N.~Seiberg,
  ``Electric--magnetic duality in supersymmetric non-Abelian gauge
  theories,''
  \textit{Nucl.~Phys.} \textbf{B435} (1995) 129
  [\texttt{hep-th/9411149}].

\bibitem{CST2011}
  C.~Csaki, Y.~Shirman, and J.~Terning,
  ``A Seiberg dual for the MSSM: partially composite $W$ and $Z$,''
  \textit{Phys.~Rev.} \textbf{D84} (2011) 095011
  [\texttt{arXiv:1106.3074}].

\end{thebibliography}

\end{document}
